% --- Archon physics formalization source begin ---
% archon:physics
% archon:covers IPhO2026Problems/problem_IPhO_2026_1_B_2.lean
% archon:source-report reports/ipho_2026/problem_IPhO_2026_1_B_2.source.json
% archon:problem-id IPhO_2026_1
% archon:part-id B.2

\chapter{Physics problem IPhO\_2026\_1\_B\_2}
\label{ch:IPhO2026Problems_problem_IPhO_2026_1_B_2}

\paragraph{Problem source.}
At one instant a positron and an electron, each of mass m and charges of equal
magnitude and opposite sign, are separated by 100*a\_0.  Their velocities are
antiparallel and perpendicular to their separation.  Each particle has angular
momentum of magnitude mu*hbar about the center of mass.  The system is isolated,
classical, non-relativistic, and has only electrostatic interaction.  The Bohr
radius is a\_0 = 4*pi*epsilon\_0*hbar\textasciicircum{}2/(m*e\textasciicircum{}2), and k = 1/(4*pi*epsilon\_0).

Current subquestion:
For mu = 15/2 the pair is unbound. Find the angle between the asymptotic relative velocity u\_infinity and the initial positron line of motion.

\paragraph{Current subquestion.}
For mu = 15/2 the pair is unbound. Find the angle between the asymptotic relative velocity u\_infinity and the initial positron line of motion.

\paragraph{Recorded answer/context.}
The signed deflection is -16.60 degrees, i.e. 16.60 degrees below the initial line of motion.

\paragraph{Figure/image path.}
/root/proposal\_for\_physic/science-mango/ipho\_2026\_source/image/T1\_page-2.png

\paragraph{Formalization target.}
create a compiling Lean file with sorry bodies at `IPhO2026Problems/problem\_IPhO\_2026\_1\_B\_2.lean`. The Lean declarations must preserve the physical quantities, dimensions or dimensional roles, figure labels, governing-law hypotheses, and final relation expressed by this problem.
Use Mathlib/Physlib names found through LeanExplore where available. If a physics API is missing, introduce faithful local abstractions rather than scalar placeholder aliases.

\begin{theorem}[Physics formalization target]
  \label{thm:physics:IPhO_2026_1_B_2:target}
  \lean{IPhO2026_1_B_2.asymptotic_relative_velocity_angle}
  \uses{def:physics:IPhO_2026_1_B_2:aux001, def:physics:IPhO_2026_1_B_2:aux002, def:physics:IPhO_2026_1_B_2:aux003, def:physics:IPhO_2026_1_B_2:aux004, def:physics:IPhO_2026_1_B_2:aux005, def:physics:IPhO_2026_1_B_2:aux006, def:physics:IPhO_2026_1_B_2:aux007, def:physics:IPhO_2026_1_B_2:aux008, def:physics:IPhO_2026_1_B_2:aux009, def:physics:IPhO_2026_1_B_2:aux010, def:physics:IPhO_2026_1_B_2:aux011, def:physics:IPhO_2026_1_B_2:aux012, def:physics:IPhO_2026_1_B_2:aux013, def:physics:IPhO_2026_1_B_2:aux014, def:physics:IPhO_2026_1_B_2:aux015, def:physics:IPhO_2026_1_B_2:aux016, def:physics:IPhO_2026_1_B_2:aux017, def:physics:IPhO_2026_1_B_2:aux018, def:physics:IPhO_2026_1_B_2:aux019, def:physics:IPhO_2026_1_B_2:aux020, def:physics:IPhO_2026_1_B_2:aux021, def:physics:IPhO_2026_1_B_2:aux022, def:physics:IPhO_2026_1_B_2:aux023, def:physics:IPhO_2026_1_B_2:aux024, def:physics:IPhO_2026_1_B_2:aux025, def:physics:IPhO_2026_1_B_2:aux026, lem:physics:IPhO_2026_1_B_2:aux027, lem:physics:IPhO_2026_1_B_2:aux028, lem:physics:IPhO_2026_1_B_2:aux029, lem:physics:IPhO_2026_1_B_2:aux030}
  For μ = 15/2, the outgoing relative velocity is deflected clockwise, below the initial positron line. The last conjunct is a rounding certificate: the signed angle in degrees rounds to -16.60° to two decimal places.
\end{theorem}
\begin{proof}
  Use the typed geometry, governing-law, branch, and measurement interfaces listed in the declaration index.  Eliminate the intermediate readouts in their physical order and then evaluate the requested exact or uncertainty relation.
\end{proof}
\section*{Formal declaration index}
The following blocks record the typed carriers and bridge declarations used by the target.  Their dependency lists follow the declaration-level model in the covered Lean file.

\begin{definition}[Declaration Plane]
  \label{def:physics:IPhO_2026_1_B_2:aux001}
  \lean{IPhO2026_1_B_2.Plane}
  The declaration Plane introduces a typed component of the physical model.
\end{definition}
\begin{proof}
  This is the typed definition of the stated physical carrier; its fields or defining equation provide the claimed data.
\end{proof}

\begin{definition}[Declaration rightward]
  \label{def:physics:IPhO_2026_1_B_2:aux002}
  \lean{IPhO2026_1_B_2.rightward}
  \uses{def:physics:IPhO_2026_1_B_2:aux001}
  The rightward direction of the initial positron velocity in Fig. 1b.
\end{definition}
\begin{proof}
  This is the typed definition of the stated physical carrier; its fields or defining equation provide the claimed data.
\end{proof}

\begin{definition}[Declaration upward]
  \label{def:physics:IPhO_2026_1_B_2:aux003}
  \lean{IPhO2026_1_B_2.upward}
  \uses{def:physics:IPhO_2026_1_B_2:aux001}
  The upward direction from the electron to the positron in Fig. 1b.
\end{definition}
\begin{proof}
  This is the typed definition of the stated physical carrier; its fields or defining equation provide the claimed data.
\end{proof}

\begin{definition}[Declaration actionDimension]
  \label{def:physics:IPhO_2026_1_B_2:aux004}
  \lean{IPhO2026_1_B_2.actionDimension}
  The dimension of action and angular momentum, M L² T⁻¹.
\end{definition}
\begin{proof}
  This is the typed definition of the stated physical carrier; its fields or defining equation provide the claimed data.
\end{proof}

\begin{definition}[Declaration energyDimension]
  \label{def:physics:IPhO_2026_1_B_2:aux005}
  \lean{IPhO2026_1_B_2.energyDimension}
  The dimension of energy, M L² T⁻².
\end{definition}
\begin{proof}
  This is the typed definition of the stated physical carrier; its fields or defining equation provide the claimed data.
\end{proof}

\begin{definition}[Declaration coulombConstantDimension]
  \label{def:physics:IPhO_2026_1_B_2:aux006}
  \lean{IPhO2026_1_B_2.coulombConstantDimension}
  The dimension of Coulomb's constant, M L³ T⁻² C⁻².
\end{definition}
\begin{proof}
  This is the typed definition of the stated physical carrier; its fields or defining equation provide the claimed data.
\end{proof}

\begin{definition}[Declaration permittivityDimension]
  \label{def:physics:IPhO_2026_1_B_2:aux007}
  \lean{IPhO2026_1_B_2.permittivityDimension}
  \uses{def:physics:IPhO_2026_1_B_2:aux006}
  The dimension of vacuum permittivity, inverse to Coulomb's constant.
\end{definition}
\begin{proof}
  This is the typed definition of the stated physical carrier; its fields or defining equation provide the claimed data.
\end{proof}

\begin{definition}[Declaration QuantityKind]
  \label{def:physics:IPhO_2026_1_B_2:aux008}
  \lean{IPhO2026_1_B_2.QuantityKind}
  An abstract physical quantity kind with dimension d, together with its numerical readout in the one coherent system of units used in the equations.
\end{definition}
\begin{proof}
  This is the typed definition of the stated physical carrier; its fields or defining equation provide the claimed data.
\end{proof}

\begin{definition}[Declaration QuantityModel]
  \label{def:physics:IPhO_2026_1_B_2:aux009}
  \lean{IPhO2026_1_B_2.QuantityModel}
  \uses{def:physics:IPhO_2026_1_B_2:aux004, def:physics:IPhO_2026_1_B_2:aux005, def:physics:IPhO_2026_1_B_2:aux006, def:physics:IPhO_2026_1_B_2:aux007, def:physics:IPhO_2026_1_B_2:aux008}
  The dimensioned primitive quantity types needed by the scattering model.
\end{definition}
\begin{proof}
  This is the typed definition of the stated physical carrier; its fields or defining equation provide the claimed data.
\end{proof}

\begin{definition}[Declaration ScatteringScenario]
  \label{def:physics:IPhO_2026_1_B_2:aux010}
  \lean{IPhO2026_1_B_2.ScatteringScenario}
  \uses{def:physics:IPhO_2026_1_B_2:aux001, def:physics:IPhO_2026_1_B_2:aux009}
  The physical objects, trajectories, conic data, and asymptotic readouts in part B.2. The initial instant is represented by t = 0, and t → +∞ describes the outgoing unbound branch.
\end{definition}
\begin{proof}
  This is the typed definition of the stated physical carrier; its fields or defining equation provide the claimed data.
\end{proof}

\begin{definition}[Declaration particleMassReadout]
  \label{def:physics:IPhO_2026_1_B_2:aux011}
  \lean{IPhO2026_1_B_2.particleMassReadout}
  \uses{def:physics:IPhO_2026_1_B_2:aux010}
  Scalar mass readout in the chosen coherent unit system.
\end{definition}
\begin{proof}
  This is the typed definition of the stated physical carrier; its fields or defining equation provide the claimed data.
\end{proof}

\begin{definition}[Declaration chargeMagnitudeReadout]
  \label{def:physics:IPhO_2026_1_B_2:aux012}
  \lean{IPhO2026_1_B_2.chargeMagnitudeReadout}
  \uses{def:physics:IPhO_2026_1_B_2:aux010}
  Positive elementary-charge magnitude readout in the chosen unit system.
\end{definition}
\begin{proof}
  This is the typed definition of the stated physical carrier; its fields or defining equation provide the claimed data.
\end{proof}

\begin{definition}[Declaration hbarReadout]
  \label{def:physics:IPhO_2026_1_B_2:aux013}
  \lean{IPhO2026_1_B_2.hbarReadout}
  \uses{def:physics:IPhO_2026_1_B_2:aux010}
  Reduced Planck constant readout.
\end{definition}
\begin{proof}
  This is the typed definition of the stated physical carrier; its fields or defining equation provide the claimed data.
\end{proof}

\begin{definition}[Declaration vacuumPermittivityReadout]
  \label{def:physics:IPhO_2026_1_B_2:aux014}
  \lean{IPhO2026_1_B_2.vacuumPermittivityReadout}
  \uses{def:physics:IPhO_2026_1_B_2:aux010}
  Vacuum-permittivity readout.
\end{definition}
\begin{proof}
  This is the typed definition of the stated physical carrier; its fields or defining equation provide the claimed data.
\end{proof}

\begin{definition}[Declaration coulombConstantReadout]
  \label{def:physics:IPhO_2026_1_B_2:aux015}
  \lean{IPhO2026_1_B_2.coulombConstantReadout}
  \uses{def:physics:IPhO_2026_1_B_2:aux010}
  Coulomb-constant readout.
\end{definition}
\begin{proof}
  This is the typed definition of the stated physical carrier; its fields or defining equation provide the claimed data.
\end{proof}

\begin{definition}[Declaration bohrRadiusReadout]
  \label{def:physics:IPhO_2026_1_B_2:aux016}
  \lean{IPhO2026_1_B_2.bohrRadiusReadout}
  \uses{def:physics:IPhO_2026_1_B_2:aux010}
  Bohr-radius readout.
\end{definition}
\begin{proof}
  This is the typed definition of the stated physical carrier; its fields or defining equation provide the claimed data.
\end{proof}

\begin{definition}[Declaration initialSeparationReadout]
  \label{def:physics:IPhO_2026_1_B_2:aux017}
  \lean{IPhO2026_1_B_2.initialSeparationReadout}
  \uses{def:physics:IPhO_2026_1_B_2:aux010}
  Initial electron--positron separation readout.
\end{definition}
\begin{proof}
  This is the typed definition of the stated physical carrier; its fields or defining equation provide the claimed data.
\end{proof}

\begin{definition}[Declaration totalAngularMomentumReadout]
  \label{def:physics:IPhO_2026_1_B_2:aux018}
  \lean{IPhO2026_1_B_2.totalAngularMomentumReadout}
  \uses{def:physics:IPhO_2026_1_B_2:aux010}
  Conserved total angular-momentum magnitude readout.
\end{definition}
\begin{proof}
  This is the typed definition of the stated physical carrier; its fields or defining equation provide the claimed data.
\end{proof}

\begin{definition}[Declaration totalEnergyReadout]
  \label{def:physics:IPhO_2026_1_B_2:aux019}
  \lean{IPhO2026_1_B_2.totalEnergyReadout}
  \uses{def:physics:IPhO_2026_1_B_2:aux010}
  Conserved total-energy readout.
\end{definition}
\begin{proof}
  This is the typed definition of the stated physical carrier; its fields or defining equation provide the claimed data.
\end{proof}

\begin{definition}[Declaration polarConicParameterReadout]
  \label{def:physics:IPhO_2026_1_B_2:aux020}
  \lean{IPhO2026_1_B_2.polarConicParameterReadout}
  \uses{def:physics:IPhO_2026_1_B_2:aux010}
  The length readout denoted by a in the supplied polar-conic hint.
\end{definition}
\begin{proof}
  This is the typed definition of the stated physical carrier; its fields or defining equation provide the claimed data.
\end{proof}

\begin{definition}[Declaration separationVector]
  \label{def:physics:IPhO_2026_1_B_2:aux021}
  \lean{IPhO2026_1_B_2.separationVector}
  \uses{def:physics:IPhO_2026_1_B_2:aux001, def:physics:IPhO_2026_1_B_2:aux010}
  Separation vector directed from the electron to the positron.
\end{definition}
\begin{proof}
  This is the typed definition of the stated physical carrier; its fields or defining equation provide the claimed data.
\end{proof}

\begin{definition}[Declaration relativeVelocity]
  \label{def:physics:IPhO_2026_1_B_2:aux022}
  \lean{IPhO2026_1_B_2.relativeVelocity}
  \uses{def:physics:IPhO_2026_1_B_2:aux001, def:physics:IPhO_2026_1_B_2:aux010}
  Positron velocity relative to the electron.
\end{definition}
\begin{proof}
  This is the typed definition of the stated physical carrier; its fields or defining equation provide the claimed data.
\end{proof}

\begin{definition}[Declaration separationRadius]
  \label{def:physics:IPhO_2026_1_B_2:aux023}
  \lean{IPhO2026_1_B_2.separationRadius}
  \uses{def:physics:IPhO_2026_1_B_2:aux010, def:physics:IPhO_2026_1_B_2:aux021}
  Electron--positron separation distance.
\end{definition}
\begin{proof}
  This is the typed definition of the stated physical carrier; its fields or defining equation provide the claimed data.
\end{proof}

\begin{definition}[Declaration signedDeflectionRadians]
  \label{def:physics:IPhO_2026_1_B_2:aux024}
  \lean{IPhO2026_1_B_2.signedDeflectionRadians}
  \uses{def:physics:IPhO_2026_1_B_2:aux010}
  The signed, counterclockwise-positive deflection from the initial positron line.
\end{definition}
\begin{proof}
  This is the typed definition of the stated physical carrier; its fields or defining equation provide the claimed data.
\end{proof}

\begin{definition}[Declaration radiansToDegrees]
  \label{def:physics:IPhO_2026_1_B_2:aux025}
  \lean{IPhO2026_1_B_2.radiansToDegrees}
  Conversion of a dimensionless angle readout from radians to degrees.
\end{definition}
\begin{proof}
  This is the typed definition of the stated physical carrier; its fields or defining equation provide the claimed data.
\end{proof}

\begin{definition}[Declaration CoulombScatteringLaws]
  \label{def:physics:IPhO_2026_1_B_2:aux026}
  \lean{IPhO2026_1_B_2.CoulombScatteringLaws}
  \uses{def:physics:IPhO_2026_1_B_2:aux002, def:physics:IPhO_2026_1_B_2:aux003, def:physics:IPhO_2026_1_B_2:aux010, def:physics:IPhO_2026_1_B_2:aux011, def:physics:IPhO_2026_1_B_2:aux012, def:physics:IPhO_2026_1_B_2:aux013, def:physics:IPhO_2026_1_B_2:aux014, def:physics:IPhO_2026_1_B_2:aux015, def:physics:IPhO_2026_1_B_2:aux016, def:physics:IPhO_2026_1_B_2:aux017, def:physics:IPhO_2026_1_B_2:aux018, def:physics:IPhO_2026_1_B_2:aux019, def:physics:IPhO_2026_1_B_2:aux020, def:physics:IPhO_2026_1_B_2:aux021, def:physics:IPhO_2026_1_B_2:aux022, def:physics:IPhO_2026_1_B_2:aux023}
  The governing physical laws and source/figure readouts. These fields state Newton--Coulomb dynamics, the initial data, conserved-energy and angular-momentum relations, the two supplied conic hints, and genuine limit relations on the outgoing branch. They do not contain the requested deflection angle.
\end{definition}
\begin{proof}
  This is the typed definition of the stated physical carrier; its fields or defining equation provide the claimed data.
\end{proof}

\begin{lemma}[Declaration outgoing\_polar\_angle\_of\_hyperbolic\_conic]
  \label{lem:physics:IPhO_2026_1_B_2:aux027}
  \lean{IPhO2026_1_B_2.outgoing_polar_angle_of_hyperbolic_conic}
  The conic equation on an unbounded branch forces the limiting polar angle to be arccos (1 / eccentricity). The range hypothesis selects the outgoing branch rather than the other root of the cosine equation.
\end{lemma}
\begin{proof}
  Substitute the cited governing relations in order and use the stated positivity, orientation, and branch conditions to obtain the conclusion.
\end{proof}

\begin{lemma}[Declaration eccentricity\_at\_mu\_fifteen\_halves]
  \label{lem:physics:IPhO_2026_1_B_2:aux028}
  \lean{IPhO2026_1_B_2.eccentricity_at_mu_fifteen_halves}
  \uses{def:physics:IPhO_2026_1_B_2:aux010, def:physics:IPhO_2026_1_B_2:aux026, lem:physics:IPhO_2026_1_B_2:aux027}
  The source constants, initial angular momentum, energy relation, and supplied eccentricity formula give eccentricity 7/2 when μ = 15/2.
\end{lemma}
\begin{proof}
  Substitute the cited governing relations in order and use the stated positivity, orientation, and branch conditions to obtain the conclusion.
\end{proof}

\begin{lemma}[Declaration outgoing\_polar\_angle\_at\_mu\_fifteen\_halves]
  \label{lem:physics:IPhO_2026_1_B_2:aux029}
  \lean{IPhO2026_1_B_2.outgoing_polar_angle_at_mu_fifteen_halves}
  \uses{def:physics:IPhO_2026_1_B_2:aux010, def:physics:IPhO_2026_1_B_2:aux026, lem:physics:IPhO_2026_1_B_2:aux027, lem:physics:IPhO_2026_1_B_2:aux028}
  The outgoing branch therefore has polar angle arccos (2/7).
\end{lemma}
\begin{proof}
  Substitute the cited governing relations in order and use the stated positivity, orientation, and branch conditions to obtain the conclusion.
\end{proof}

\begin{lemma}[Declaration fig1b\_signed\_deflection\_from\_polar\_angle]
  \label{lem:physics:IPhO_2026_1_B_2:aux030}
  \lean{IPhO2026_1_B_2.fig1b_signed_deflection_from_polar_angle}
  \uses{def:physics:IPhO_2026_1_B_2:aux010, def:physics:IPhO_2026_1_B_2:aux024, def:physics:IPhO_2026_1_B_2:aux026, lem:physics:IPhO_2026_1_B_2:aux027, lem:physics:IPhO_2026_1_B_2:aux028, lem:physics:IPhO_2026_1_B_2:aux029}
  Fig. 1b puts the initial positron velocity along the positive horizontal axis, while the conic's polar zero-axis points downward. The radial outgoing limit therefore converts polar angle θ∞ into signed deflection θ∞ - π/2.
\end{lemma}
\begin{proof}
  Substitute the cited governing relations in order and use the stated positivity, orientation, and branch conditions to obtain the conclusion.
\end{proof}
% --- Archon physics formalization source end ---
